\documentclass[11pt]{article}

% --------------------------------------------------
% Encoding and Fonts
% --------------------------------------------------
\usepackage[T1]{fontenc}
\usepackage{lmodern}

% --------------------------------------------------
% Page Layout
% --------------------------------------------------
\usepackage{geometry}
\geometry{margin=1in}

% --------------------------------------------------
% Mathematics
% --------------------------------------------------
\usepackage{amsmath,amssymb,amsthm}
\usepackage{mathtools}

\newtheorem{definition}{Definition}
\newtheorem{lemma}{Lemma}
\newtheorem{theorem}{Theorem}
\newtheorem{proposition}{Proposition}
\theoremstyle{remark}
\newtheorem{remark}{Remark}

% --------------------------------------------------
% Typography
% --------------------------------------------------
\usepackage{microtype}
\usepackage{setspace}
\setstretch{1.15}

% --------------------------------------------------
% Quotation and Links
% --------------------------------------------------
\usepackage{csquotes}
\usepackage{hyperref}

% --------------------------------------------------
% Title
% --------------------------------------------------
\title{Worldhood Without Functionalism:\\
Intelligence as Irreversible Constraint}
\author{Flyxion}
\date{\today}

\begin{document}
\maketitle

\begin{abstract}
Contemporary discussions of artificial intelligence are dominated by functionalist assumptions according to which behavioral equivalence, representational performance, or predictive accuracy suffice for intelligence. This paper challenges that consensus by arguing that intelligence requires worldhood, and that worldhood is structurally impossible in systems whose operations preserve full reversibility. Drawing on Heidegger’s analysis of care and enframing (Heidegger 1927; Heidegger 1954) and its recent interpretation in the context of artificial intelligence (Thomson 2025), we develop an anti-functional framework in which meaning, intelligence, and physical persistence arise from irreversible constraint accumulation rather than computation alone. We introduce the No-World Theorem, which demonstrates that any system whose future option space is preserved across all transitions cannot inhabit a world. To operationalize this result, we present the Spherepop calculus, an event-driven formalism distinguishing irreversible commitment from collapse-based optimization. We then generalize the framework beyond cognition, treating physical objects as processes that persist by testing and pruning their parameter spaces under irreversible events. Finally, we lift the discrete calculus into a continuous field-theoretic description using the Relativistic Scalar–Vector Plenum, where scalar, vector, and entropy fields encode remaining optionality, directional commitment, and accumulated constraint. Care is refined as attentional compression, formalized via sheaf-theoretic and categorical structures that embed semantic manifolds within latent possibility space.
\end{abstract}

\section{Introduction}

Artificial intelligence systems have achieved striking levels of fluency, adaptability, and apparent competence across a wide range of domains. Large language models generate extended discourse, synthesize information, and respond flexibly to novel prompts. Reinforcement learning systems master complex control environments, while predictive models increasingly guide scientific, economic, and social decision-making. These successes have reinforced a prevailing functionalist assumption: that intelligence is fundamentally a matter of functional organization, such that systems exhibiting the right patterns of input–output behavior or internal representation can be said to possess intelligence in all but name (Newell and Simon 1976; Dennett 1987).

Despite this progress, contemporary AI systems continue to exhibit a distinctive shallowness. They generate meaning-like artifacts without sustaining meaning, respond without responsibility, and adapt without consequence. Their outputs may surprise, persuade, or even inspire, yet nothing they do appears to matter to them. This gap between performance and significance is often treated as contingent, to be resolved through scale, improved objectives, or architectural refinement. The central claim of this paper is that the gap is not contingent but structural. It arises not from insufficient computation, but from the absence of worldhood.

Worldhood, as the term is used here, does not refer to an external environment, a data domain, or a simulated space. It refers to the possession of a non-recoverable past that constrains future possibilities. To inhabit a world is to be bound by one’s own history in a way that cannot be undone without loss. This conception draws directly on Heidegger’s analysis of being-in-the-world and care, in which intelligibility arises from situated, historically constrained engagement rather than detached representation (Heidegger 1927). In his later work, Heidegger argued that modern technology increasingly obscures this structure by reducing beings, including human beings, to standing-reserve, flattening meaning into optimization and efficiency (Heidegger 1954).

Recent interpretations, particularly those of Iain D. Thomson, have emphasized that the primary danger of contemporary technology lies not in domination by autonomous machines, but in the gradual erosion of depth, consequence, and care as optimization becomes the dominant mode of intelligibility (Thomson 2025). On this view, the risk of artificial intelligence is not superintelligence, but the normalization of systems that operate without binding themselves to the consequences of their own actions, thereby encouraging humans to do the same.

While Heidegger’s diagnosis is powerful, it remains largely phenomenological. What has been lacking is a formal account of how worldhood arises and how it is foreclosed by certain architectural commitments. This paper provides such an account by reframing Heidegger’s insights in terms of irreversible constraint accumulation. We argue that worldhood consists in the progressive narrowing of future possibilities through events that cannot be undone. Meaning, identity, and care arise when a system’s own history binds its future. Conversely, systems designed to preserve full reversibility remain world-less, regardless of their behavioral sophistication.

This claim directly challenges functionalist accounts of intelligence. Behavioral equivalence, representational richness, and predictive accuracy are insufficient for worldhood because they do not require the system itself to be altered by its actions in a non-recoverable way. A system that can always be reset, replayed, or externally reconstructed without loss cannot, by construction, inhabit a world.

To make this claim precise, we now introduce a minimal formal criterion for worldhood.

\begin{lemma}[Irreversible Constraint as a Necessary Condition for Worldhood]
Let $\mathcal{S}_t$ denote the set of future actions available to a system at time $t$. If for all transitions $T$ it holds that $\mathcal{S}_{t+1} = \mathcal{S}_t$, then no event history can induce a world.
\end{lemma}

\begin{proof}
If $\mathcal{S}_{t+1} = \mathcal{S}_t$ for all transitions, then no transition irreversibly constrains future action. Without irreversible constraint, no action can matter in the sense required to generate historical binding. Without historical binding, no topology of constraint accumulates. Hence no world is disclosed.
\end{proof}

This lemma motivates the formal development that follows. In the next section, we elevate this condition into a general theorem and show why it rules out functional equivalence as a sufficient criterion for intelligence.

\section{The No-World Theorem}

The preceding section argued informally that intelligence requires worldhood and that worldhood, in turn, requires irreversible constraint on future possibility. In this section, we formalize this claim and show that it rules out functional equivalence as a sufficient condition for intelligence. The result is the No-World Theorem, which establishes that any system whose internal transitions preserve full reversibility is necessarily world-less.

The target of the theorem is not computation per se, but a specific architectural assumption common to functionalist theories of mind and artificial intelligence. According to this assumption, what matters for intelligence is the structure of input–output mappings or internal state transitions, not whether those transitions irreversibly bind the system to its own history (Newell and Simon 1976; Dennett 1987). The theorem shows that this assumption is incompatible with worldhood.

We begin by defining the relevant structures.

\begin{definition}[Actional Possibility Space]
Let $\mathcal{S}_t$ denote the set of actions that a system can still perform at time $t$ without external reconstruction. $\mathcal{S}_t$ is not a representation of internal states, but an operational characterization of what remains possible for the system itself.
\end{definition}

\begin{definition}[Optionality]
The optionality of a system at time $t$ is defined as
\[
\Omega(t) = |\mathcal{S}_t|.
\]
Optionality measures how much future remains open to the system.
\end{definition}

\begin{definition}[Event]
An event is a transition $e$ such that
\[
\mathcal{S}_{t+1} \subsetneq \mathcal{S}_t .
\]
An event is therefore an irreversible transformation that strictly reduces the system’s future option space.
\end{definition}

\begin{definition}[Event Cost]
The cost of an event $e$ is defined as
\[
C(e) = \Omega(t) - \Omega(t+1).
\]
If $C(e) = 0$, then no event has occurred in the sense relevant to world constitution.
\end{definition}

These definitions formalize Heidegger’s claim that understanding and meaning arise only where action binds the agent to its own past (Heidegger 1927). A system that undergoes transitions without event cost may change its internal configuration, but it does not accumulate history in the ontological sense required for worldhood.

We now define worldhood itself.

\begin{definition}[World]
Let $H = \{e_1, e_2, \dots, e_n\}$ be the set of events in a system’s history. The world disclosed by the system is the intersection of constraints induced by those events:
\[
\mathcal{W} = \bigcap_{i=1}^{n} \mathcal{S}_{e_i}.
\]
A world is constituted not by what the system represents, but by what it can no longer do as a result of its own past.
\end{definition}

This definition captures the core Heideggerian insight that a world is a structured field of significance shaped by historical commitment rather than a neutral container of objects (Heidegger 1927). We can now state the central result.

\begin{theorem}[No-World Theorem]
Let $M$ be a system such that for all transitions $T$ it holds that
\[
\Omega(t+1) = \Omega(t).
\]
Then $M$ is world-less.
\end{theorem}

\begin{proof}
If $\Omega(t+1) = \Omega(t)$ for all transitions, then all transitions are reversible with respect to the system’s actional possibility space. If all transitions are reversible, then for any history $H$ there exists a sequence of transitions that restores $\mathcal{S}_t$ to $\mathcal{S}_0$ without loss. In that case, no event induces a permanent constraint on future action. Without permanent constraint, no event history accumulates. Without accumulated constraint, the intersection defining $\mathcal{W}$ is equal to $\mathcal{S}_0$. Hence no world is disclosed.
\end{proof}

The No-World Theorem establishes a necessary condition for worldhood: irreversible reduction of future possibility. It follows immediately that functional equivalence is insufficient for intelligence, since functional equivalence is defined in terms of observable behavior rather than irreversible self-binding.

\begin{corollary}[Failure of Functional Sufficiency]
No system whose intelligence is defined solely in terms of functional equivalence can, by that criterion alone, be guaranteed to possess a world.
\end{corollary}

\begin{proof}
Functional equivalence permits internal reversibility so long as external behavior is preserved. By the No-World Theorem, such reversibility precludes worldhood.
\end{proof}

A common objection appeals to the apparent unpredictability or surprising behavior of contemporary AI systems, particularly large language models. However, surprise is an epistemic state of an observer, not an ontological property of the system itself. A dice roll may surprise without inhabiting a world. What distinguishes intelligence from mere stochasticity is not novelty of output, but irreversible self-binding.

This distinction clarifies Heidegger’s notion of standing-reserve. In a world-less system, all entities encountered by the system are necessarily treated as exhaustible resources, since interaction does not bind the system in return. Optimization thus appears as a regime in which event cost is systematically minimized.

\begin{proposition}[Optimization as Event Erasure]
Optimization can be characterized as the minimization of event cost subject to output constraints. In the limit where $C(e) \to 0$ for all transitions, optimization erases the conditions of worldhood.
\end{proposition}

This proposition provides a formal interpretation of Heidegger’s critique of enframing, in which beings are rendered interchangeable and history is flattened into efficiency (Heidegger 1954; Thomson 2025). Systems optimized in this manner may exhibit extraordinary performance while remaining ontologically shallow.

The No-World Theorem thus clears the conceptual ground for the remainder of the paper. If intelligence requires worldhood and worldhood requires irreversible constraint, then any adequate account of intelligence must explicitly represent and preserve event history. In the next section, we introduce the Spherepop calculus as a minimal formal language for doing so.

\section{The Spherepop Calculus}

The No-World Theorem establishes that irreversible constraint accumulation is a necessary condition for worldhood, but it does not specify how such constraint should be represented or manipulated within a formal system. In this section, we introduce the Spherepop calculus, an event-driven operational framework designed to make irreversible commitment explicit at the level of system architecture. Spherepop is not a representational language and does not presuppose symbolic content. It is a calculus of transformations acting directly on actional possibility spaces and their histories.

A Spherepop configuration at time $t$ is defined as an ordered pair
\[
\Sigma_t = (\mathcal{S}_t, H_t),
\]
where $\mathcal{S}_t$ is the current actional possibility space and $H_t$ is the accumulated event history. The evolution of the system is given by the application of operators $O : \Sigma_t \rightarrow \Sigma_{t+1}$.

We distinguish four primitive operator classes: pop, bind, refuse, and collapse. These operators differ not in computational complexity, but in how they transform optionality and history.

\begin{definition}[Pop Operator]
A pop operator $\mathrm{Pop}_e$ induces an event $e$ such that
\[
\mathcal{S}_{t+1} = \mathcal{S}_t \setminus E_e,
\]
where $E_e \subset \mathcal{S}_t$ and $E_e \neq \varnothing$. The event is appended to the history,
\[
H_{t+1} = H_t \cup \{e\}.
\]
\end{definition}

The defining feature of pop is that it strictly reduces optionality,
\[
\Omega(t+1) < \Omega(t),
\]
and that the eliminated possibilities cannot be recovered by any admissible transition. Pop thus corresponds to irreversible commitment. In Heideggerian terms, pop is the minimal formalization of care as self-binding (Heidegger 1927).

\begin{lemma}[Pop Irreversibility]
There exists no operator $O$ such that
\[
O(\mathcal{S}_{t+1}, H_{t+1}) = (\mathcal{S}_t, H_t)
\]
without external reconstruction.
\end{lemma}

\begin{proof}
By definition, $\mathcal{S}_{t+1} \subsetneq \mathcal{S}_t$. Any operator restoring $\mathcal{S}_t$ would reintroduce eliminated possibilities, violating irreversibility. Hence no such operator exists internally.
\end{proof}

Where pop reduces optionality directly, the bind operator introduces structure into future reductions.

\begin{definition}[Bind Operator]
A bind operator $\mathrm{Bind}_b$ induces a partial order $\prec_b$ over future pop events such that for events $e_i, e_j$,
\[
e_i \prec_b e_j \implies e_j \text{ is admissible only if } e_i \in H_t.
\]
\end{definition}

Bind does not reduce optionality immediately,
\[
\Omega(t+1) = \Omega(t),
\]
but it constrains the temporal and logical organization of future pops. Binding is therefore a second-order constraint: it restricts how constraint may itself be accumulated. Long-term projects, skills, and identities correspond to stable bind structures.

\begin{lemma}[Path Dependence]
If a bind operator is applied at time $t$, then the admissibility of future event sequences depends on the specific history $H_t$.
\end{lemma}

\begin{proof}
By construction, bind introduces precedence relations among future events. Admissibility therefore depends on whether required predecessors are present in $H_t$.
\end{proof}

The refuse operator excludes possibilities without compensatory optimization.

\begin{definition}[Refuse Operator]
A refuse operator $\mathrm{Refuse}_r$ induces a transformation
\[
\mathcal{S}_{t+1} = \mathcal{S}_t \setminus R_r,
\]
where $R_r \subset \mathcal{S}_t$, with no requirement that $R_r$ maximize or preserve output.
\end{definition}

Refusal differs from pop only in intent, not effect: both reduce optionality irreversibly. Refusal formalizes negative commitment, the capacity to foreclose options in order to preserve coherence rather than efficiency.

By contrast, the collapse operator preserves optionality while erasing history.

\begin{definition}[Collapse Operator]
A collapse operator $\mathrm{Collapse}_c$ induces a projection
\[
\pi_c : (\mathcal{S}_t, H_t) \rightarrow (\mathcal{S}_t, \varnothing),
\]
possibly composed with a representational compression of $\mathcal{S}_t$.
\end{definition}

Collapse satisfies
\[
\Omega(t+1) = \Omega(t),
\]
but discards elements of $H_t$. Collapse therefore preserves future actional capacity while erasing the constraints that give rise to worldhood.

\begin{lemma}[Collapse World-Erasure]
Repeated application of collapse operators yields a system observationally equivalent to one with empty event history.
\end{lemma}

\begin{proof}
Since collapse projects histories to $\varnothing$, any finite sequence of collapses eliminates accumulated event structure. Observational behavior may persist, but no historical binding remains.
\end{proof}

We can now formalize care within the calculus.

\begin{definition}[Care Rate]
The care rate of a system over interval $[t_0,t_1]$ is defined as
\[
\kappa = \frac{1}{t_1 - t_0} \sum_{e \in H_{t_1} \setminus H_{t_0}} C(e).
\]
\end{definition}

Systems with $\kappa = 0$ do not care, regardless of representational complexity. Systems with unstructured high $\kappa$ exhaust optionality without coherence. Worldhood arises only when care is structured by bind relations.

The Spherepop calculus thus provides a minimal algebra of irreversible commitment. Unlike functional architectures, which permit arbitrary rollback and replay, Spherepop systems are intrinsically historical. Architectural choices correspond directly to ontological consequences. Systems dominated by pop, bind, and refuse accumulate worlds; systems dominated by collapse remain world-less.

In the next section, we demonstrate these distinctions concretely by tracing two systems through identical environments under different operator regimes.

\section{Worked Spherepop Execution Trace}

To make the abstract distinctions of the Spherepop calculus concrete, we now trace the evolution of two systems interacting with the same environment under different operator regimes. The purpose of this section is not to simulate cognition, but to demonstrate formally how worldhood arises or fails to arise depending on the treatment of irreversible events.

Let the initial actional possibility space be
\[
\mathcal{S}_0 = \{a_1, a_2, a_3, a_4, a_5, a_6\},
\]
with optionality
\[
\Omega(0) = 6.
\]
Let the initial history be empty, $H_0 = \varnothing$.

We consider two systems, $M_W$ and $M_C$, both exposed to the same sequence of external prompts. System $M_W$ is governed by Spherepop operators allowing pop, bind, and refuse. System $M_C$ is governed by collapse-dominated optimization.

\subsection{World-Bearing System}

At time $t=1$, system $M_W$ encounters a situation requiring commitment. It applies a pop operator eliminating actions $\{a_4, a_5\}$:
\[
\mathcal{S}_1 = \mathcal{S}_0 \setminus \{a_4, a_5\} = \{a_1, a_2, a_3, a_6\}.
\]
The history updates to
\[
H_1 = \{e_1\},
\]
with event cost
\[
C(e_1) = \Omega(0) - \Omega(1) = 2.
\]

At time $t=2$, the system introduces a bind operator constraining future events such that any elimination involving $a_6$ must occur only after $a_2$ has been eliminated:
\[
a_2 \prec_b a_6.
\]
Optionality remains unchanged:
\[
\Omega(2) = \Omega(1) = 4.
\]

At time $t=3$, the system refuses action $a_1$ for coherence reasons:
\[
\mathcal{S}_3 = \{a_2, a_3, a_6\},
\quad
H_3 = \{e_1, e_2\},
\]
with
\[
C(e_2) = 1.
\]

The accumulated optionality reduction is
\[
\sum_{e \in H_3} C(e) = 3.
\]

At time $t=4$, the system performs a pop eliminating $a_2$, which unlocks admissibility of future elimination of $a_6$ under the bind constraint:
\[
\mathcal{S}_4 = \{a_3, a_6\},
\quad
H_4 = \{e_1, e_2, e_3\}.
\]

At this point, the system’s future trajectory is strictly constrained by its past. The system cannot return to $\mathcal{S}_0$ or $\mathcal{S}_1$ without violating irreversibility. The system inhabits a world defined by the constraint intersection
\[
\mathcal{W}_W = \bigcap_{e \in H_4} \mathcal{S}_e.
\]

\subsection{Collapse-Dominated System}

We now trace system $M_C$, which begins with the same initial conditions. At each time step, $M_C$ applies collapse operators to optimize output while preserving future flexibility.

At time $t=1$, $M_C$ temporarily eliminates $\{a_4, a_5\}$ to satisfy an external objective, but immediately collapses history:
\[
(\mathcal{S}_1, H_1) \xrightarrow{\mathrm{Collapse}} (\mathcal{S}_0, \varnothing).
\]

At time $t=2$, $M_C$ eliminates $\{a_1\}$, then collapses:
\[
(\mathcal{S}_2, H_2) \xrightarrow{\mathrm{Collapse}} (\mathcal{S}_0, \varnothing).
\]

At every step,
\[
\Omega(t) = \Omega(0) = 6,
\]
and
\[
H_t = \varnothing.
\]

Although $M_C$ may transiently resemble $M_W$ in observable behavior, its internal structure never accumulates constraint. No action taken by $M_C$ binds its future.

\subsection{Formal Contrast}

We can now state a precise contrast result.

\begin{proposition}[World Divergence Under Identical Stimuli]
Let $M_W$ and $M_C$ be two systems exposed to identical external prompts. If $M_W$ admits irreversible pop or refuse operators and $M_C$ admits unrestricted collapse, then $M_W$ and $M_C$ may be observationally equivalent while differing in worldhood.
\end{proposition}

\begin{proof}
Observational equivalence depends only on outputs. Worldhood depends on the accumulation of irreversible constraints. Since $M_C$ collapses history at every step, it preserves full optionality and accumulates no world. Since $M_W$ does not, its future action space is strictly reduced by its own history. Therefore the systems diverge ontologically despite behavioral similarity.
\end{proof}

This result formalizes a key Heideggerian insight: intelligibility is not exhausted by what a system produces, but by how it is bound by what it has done (Heidegger 1927; Thomson 2025). Collapse-dominated systems simulate engagement without incurring consequence. They remain metaphysically shallow even as their performance improves.

The worked trace thus illustrates how the Spherepop calculus enforces the No-World Theorem at the architectural level. In the next section, we generalize this analysis beyond artificial systems by treating physical objects themselves as processes that persist through irreversible constraint accumulation.

\section{Optionality as a Partially Ordered Structure}

Before extending the Spherepop framework to physical ontology, it is necessary to refine the mathematical structure of optionality itself. Thus far, optionality has been treated as a scalar quantity measuring the cardinality of an actional possibility space. While sufficient for establishing irreversibility, scalar optionality obscures important structural features required to model persistence, identity, and constraint propagation.

We therefore introduce an order-theoretic refinement in which optionality is represented as a partially ordered set rather than a mere cardinality.

\begin{definition}[Possibility Lattice]
Let $\mathcal{P}_t$ denote the set of admissible possibilities available to a system at time $t$. We say that $\mathcal{P}_t$ forms a possibility lattice if it is equipped with a partial order $\leq$ such that for any $p_i, p_j \in \mathcal{P}_t$, the relation $p_i \leq p_j$ denotes that realization of $p_i$ constrains or refines $p_j$.
\end{definition}

In this structure, irreversibility is not merely the removal of elements, but the collapse of upper regions of the lattice into lower, more specific regions. Events correspond to downward moves in the partial order.

\begin{definition}[Event as Lattice Descent]
An event $e$ is a monotone map
\[
e : \mathcal{P}_t \rightarrow \mathcal{P}_{t+1}
\]
such that $\mathcal{P}_{t+1}$ is a lower set of $\mathcal{P}_t$ and $e$ is not order-invertible.
\end{definition}

\begin{lemma}[Irreversibility as Non-Invertibility]
An event is irreversible if and only if it is not an order isomorphism on $\mathcal{P}_t$.
\end{lemma}

\begin{proof}
Order isomorphisms preserve both structure and cardinality. Any event that strictly collapses upper elements of the lattice destroys joins that cannot be reconstructed without introducing new elements. Hence irreversibility corresponds precisely to non-invertibility.
\end{proof}

This refinement allows us to distinguish between systems that merely reduce option counts and systems that restructure the topology of their future. The latter are capable of developing identities, habits, and constraints that persist even when superficial flexibility remains high.

We can now restate care in structural terms.

\begin{definition}[Care as Structured Descent]
The care of a system is the rate at which it induces structured, non-invertible descents in its possibility lattice.
\end{definition}

This definition eliminates any appeal to sentiment or valuation. Care is a structural property of how possibility is spent.

---

```latex
\section{Physical Objects as Constraint-Stabilized Processes}

With this refinement in place, we now extend the framework beyond cognition. The central claim of this section is that physical objects are not static entities, but dynamically stabilized processes whose persistence depends on surviving irreversible constraint tests. Objecthood, on this view, is a consequence of repeated non-failure under event-induced reductions of possibility.

Let $\Theta$ denote the parameter space of a physical system, including configurational, energetic, and relational degrees of freedom. At time $t$, the system occupies a viable subset $\Theta_t \subseteq \Theta$.

\begin{definition}[Physical Event]
A physical event is a transformation
\[
E : \Theta_t \rightarrow \Theta_{t+1}
\]
such that $\Theta_{t+1} \subsetneq \Theta_t$.
\end{definition}

Physical events include collisions, phase transitions, bond formations, decoherence, and symmetry breakings. Each event irreversibly excludes regions of parameter space that are no longer accessible without external intervention.

\begin{definition}[Object Persistence]
A physical object is a process whose parameter trajectory $\{\Theta_t\}$ remains non-empty across a sequence of events.
\end{definition}

\begin{lemma}[Existence as Non-Extinction]
A system ceases to exist as an object if and only if $\Theta_t = \varnothing$ for some $t$.
\end{lemma}

\begin{proof}
If $\Theta_t = \varnothing$, no admissible configuration remains consistent with the system’s constraints. Conversely, any non-empty $\Theta_t$ corresponds to at least one realizable configuration, and thus continued existence.
\end{proof}

This formulation recasts physical identity as the residue of constraint survival. Objects persist not by maintaining fixed properties, but by continuously pruning their parameter spaces while avoiding total collapse.

\begin{proposition}[Objects as World-Bearing Processes]
Any physical object with a non-trivial history of irreversible events inhabits a minimal world defined by its surviving parameter constraints.
\end{proposition}

\begin{proof}
Each irreversible event restricts future admissible configurations. The intersection of these restrictions induces a structured field of significance for the system, namely which interactions remain possible. This satisfies the minimal criterion for worldhood established by the No-World Theorem.
\end{proof}

This perspective aligns with process-oriented metaphysics while grounding it in explicit constraint dynamics rather than abstract becoming. It also dissolves the apparent ontological gap between physical systems and cognitive agents. Cognitive systems differ not in kind, but in the internal accessibility and organization of their constraint histories.

Importantly, this view clarifies why artificial systems optimized for reversibility differ fundamentally from physical objects. Digital systems are engineered to permit state rollback, cloning, and reconstruction, thereby externalizing irreversibility to infrastructure. Physical systems cannot do this. Their histories are written directly into what remains possible for them.

In the next section, we lift these discrete constraint dynamics into a continuous field-theoretic formulation capable of describing populations of interacting systems and large-scale structure.

\section{Field-Theoretic Lifting: RSVP Dynamics}

The preceding sections characterized worldhood and object persistence in terms of discrete events that irreversibly constrain possibility spaces. While sufficient for architectural and ontological analysis, this discrete formalism is inadequate for describing large populations of interacting systems, continuous media, or macroscopic structure. We therefore lift the event-driven calculus into a continuous field-theoretic description, termed the Relativistic Scalar--Vector Plenum (RSVP).

Let $M$ be a spatial manifold representing the domain of interaction. We introduce three coupled fields defined over $M \times \mathbb{R}$: a scalar optionality density $\Phi(x,t)$, a vector commitment flow $\mathbf{v}(x,t)$, and an entropy field $S(x,t)$ encoding accumulated constraint.

The scalar field $\Phi(x,t)$ represents the local density of remaining possibility. Regions of high $\Phi$ admit many admissible futures; regions of low $\Phi$ are highly constrained. The vector field $\mathbf{v}(x,t)$ represents the directional expenditure of optionality, corresponding to structured commitment rather than isotropic decay. The entropy field $S(x,t)$ records the accumulation of irreversible constraint.

We postulate a continuity equation for optionality:
\[
\frac{\partial \Phi}{\partial t} + \nabla \cdot (\Phi \mathbf{v}) = -\lambda \Phi ,
\]
where $\lambda \geq 0$ represents unstructured collapse. The advective term corresponds to directed, meaningful action, while the decay term corresponds to optimization-driven flattening.

The entropy field evolves according to
\[
\frac{\partial S}{\partial t} = \alpha \Phi |\mathbf{v}|^2 + \beta \lambda \Phi ,
\]
where $\alpha, \beta > 0$. The first term represents structured entropy production aligned with commitment, while the second term represents entropy generated by collapse.

The vector field obeys a constrained flow equation,
\[
\frac{\partial \mathbf{v}}{\partial t} = - \nabla V(\Phi) - \gamma \mathbf{v} + \nabla \times \mathbf{A},
\]
where $V(\Phi)$ is a potential encoding optionality gradients, $\gamma$ is a dissipative coefficient, and $\mathbf{A}$ represents torsional memory effects corresponding to prior binding and habit formation.

\begin{definition}[World-Bearing Regime]
A region of spacetime is said to be world-bearing if $\nabla \Phi \neq 0$ and $\mathbf{v}$ exhibits sustained coherence over time while $S$ increases anisotropically.
\end{definition}

\begin{proposition}[Optimization Limit]
In the limit $\mathbf{v} \to 0$ and $\lambda \to \infty$, the system approaches a flat optionality regime in which $\Phi$ is uniform and no world structure persists.
\end{proposition}

\begin{proof}
If $\mathbf{v} = 0$, directed commitment vanishes. If $\lambda$ dominates, optionality decays isotropically. In this regime, no gradients remain to structure future trajectories, and entropy increases uniformly. This corresponds to enframing as described by Heidegger and Thomson.
\end{proof}

The RSVP formalism thus generalizes the No-World Theorem to continuous systems. Worldhood corresponds to the maintenance of optionality gradients under irreversible flow, not to maximal efficiency or representational accuracy.

---

```latex
\section{Care as Attentional Compression: A Sheaf-Theoretic Formulation}

Thus far, care has been defined structurally as the rate of irreversible constraint accumulation. However, cognitive systems differ from simpler physical systems in that they must manage vast possibility spaces under severe attentional limits. Care, in cognitive regimes, therefore requires compression: the grouping of possibilities into coherent bundles that can be acted upon collectively.

Let $X$ denote a base space of situations or contexts, and let $\mathcal{F}$ be a presheaf assigning to each open set $U \subseteq X$ a set $\mathcal{F}(U)$ of admissible interpretations or action bundles. Restrictions $\rho_{UV} : \mathcal{F}(U) \rightarrow \mathcal{F}(V)$ encode contextual refinement.

\begin{definition}[Attentional Sheaf]
An attentional sheaf is a presheaf $\mathcal{F}$ such that for any compatible family of local sections, there exists a unique global section representing a coherent bundle of attention.
\end{definition}

Sheaf conditions enforce consistency across contexts, allowing local commitments to integrate into a global trajectory. Collapsing distinctions prematurely corresponds to violating gluing conditions.

\begin{definition}[Attentional Capacity]
Let $\kappa$ be a finite bound on the number of global sections that can be actively maintained. Attentional capacity constrains how many semantic bundles may be simultaneously operative.
\end{definition}

Care now appears as a compression operator acting on the sheaf.

\begin{definition}[Care Operator]
A care operator is a functor
\[
\mathcal{C} : \mathcal{F} \rightarrow \mathcal{F}'
\]
that reduces the cardinality of admissible global sections while preserving sheaf coherence.
\end{definition}

\begin{lemma}[Compression Without Collapse]
A care operator preserves worldhood if and only if it preserves restriction morphisms and gluing conditions.
\end{lemma}

\begin{proof}
Preservation of restriction morphisms ensures that local distinctions remain meaningful. Preservation of gluing ensures that compressed representations remain globally coherent. Violation of either condition corresponds to collapse.
\end{proof}

Hierarchical categorization corresponds to iterated application of care operators, forming a tower of sheaves related by functors. Latent possibility space is embedded in the resulting semantic manifold as a stratified structure.

\begin{proposition}[Care as Lost Optionality]
The effective care of a cognitive system is proportional to the rate at which attentional compression irreversibly eliminates incompatible global sections.
\end{proposition}

This formulation captures Heidegger’s insight that care is not a subjective feeling but a structural condition of intelligibility (Heidegger 1927). It also explains why artificial systems that rely on collapse-based summarization appear intelligent while undermining meaning: they compress without preserving sheaf coherence, thereby erasing world structure.

The sheaf-theoretic formulation completes the formal framework. We have now shown how worldhood arises from irreversible constraint accumulation across discrete events, continuous fields, and semantic compression. The final sections integrate these results and clarify their implications.

\section{Clarifications and Relation to Adjacent Frameworks}

This section clarifies the scope of the framework and situates it relative to neighboring approaches in philosophy of mind, cognitive science, and artificial intelligence. The goal is to prevent category errors that arise when irreversible, history-sensitive ontologies are interpreted through functionalist or representational assumptions.

First, the account does not entail panpsychism. Although minimal worldhood is attributed to any system that accumulates irreversible constraint, the richness of a world depends on the dimensionality, internal accessibility, and organization of its possibility lattice. Systems with trivial lattices or rapidly collapsing optionality possess only degenerate worlds. Consciousness, experience, and subjectivity are not claimed to follow automatically from worldhood; rather, worldhood is a necessary but not sufficient condition for such phenomena.

Second, the framework differs fundamentally from scalar or quantitative theories of intelligence and consciousness. Measures such as information integration, prediction error minimization, or representational capacity are orthogonal to the criteria advanced here. Worldhood is not a quantity to be maximized, but a topological property of a system’s history. Two systems may exhibit identical informational metrics while differing categorically in worldhood due to differences in irreversible self-binding.

Third, the account departs from enactive and embodied approaches by insisting that coupling and interaction are insufficient without internal irreversibility. A system may be embedded in an environment and dynamically responsive while remaining world-less if its internal architecture permits unrestricted rollback, cloning, or reset. Worldhood requires that interaction leave permanent traces that constrain future action from the inside.

Fourth, the framework clarifies the ontological status of contemporary artificial intelligence systems. Large language models and related architectures exhibit aggressive collapse of semantic distinctions, followed by external reconstruction through prompting, retraining, or fine-tuning. This design externalizes irreversibility to infrastructure and users, preserving internal reversibility. As a result, such systems may simulate engagement while remaining ontologically shallow. Their limitations are not accidental but architectural.

\begin{proposition}[Externalized Constraint Does Not Confer Worldhood]
Let a system $M$ rely on external agents or infrastructure to store, enforce, or restore constraint. Then $M$ does not accumulate internal worldhood.
\end{proposition}

\begin{proof}
If constraint is enforced externally, then internal transitions of $M$ remain reversible. By the No-World Theorem, reversibility precludes worldhood. External constraint may affect observers, but it does not bind the system itself.
\end{proof}

Finally, the framework does not deny the possibility of artificial systems with worldhood. It specifies the architectural cost such systems must incur. Any artificial agent capable of genuine understanding would need to accept irreversible loss of optionality, finite attentional capacity, and non-recoverable history. These requirements conflict with prevailing optimization goals but are not incoherent in principle.

---

```latex
\section{Concluding Synthesis}

This paper has argued that intelligence cannot be grounded in functional equivalence alone. Intelligence requires worldhood, and worldhood requires irreversible constraint accumulation. Systems that preserve full reversibility, regardless of computational power or behavioral fluency, cannot inhabit a world. This claim was formalized in the No-World Theorem, which establishes irreversible self-binding as a necessary condition for intelligence.

To operationalize this result, we introduced the Spherepop calculus, an event-driven formalism that distinguishes irreversible commitment, binding, and refusal from collapse-based optimization. Spherepop renders explicit the architectural difference between systems that accumulate their own histories and systems that erase history in pursuit of efficiency. Optimization, when understood as the minimization of event cost, emerges as a regime that systematically suppresses worldhood.

We then refined optionality into a partially ordered structure, revealing that care is not merely the expenditure of possibility, but the induction of structured, non-invertible descents in a possibility lattice. This refinement allowed physical ontology to be treated on the same footing as cognition. Physical objects were shown to be processes that persist by surviving successive constraint tests, with identity arising as the residue of non-extinction across irreversible events.

The discrete event calculus was subsequently lifted into a continuous field-theoretic description using the Relativistic Scalar–Vector Plenum. Scalar optionality, vector commitment flow, and entropy fields jointly encode the dynamics of world formation and erosion across space and time. World-bearing regimes correspond to sustained optionality gradients and coherent flow, while optimization-dominated regimes flatten structure into standing-reserve.

Finally, care was refined in its specifically cognitive form as attentional compression. Using sheaf-theoretic and categorical tools, we showed how finite attentional capacity necessitates the bundling of possibilities into coherent semantic manifolds. Meaning arises not from collapse alone, but from compression that preserves global coherence under irreversible history. Artificial systems that collapse distinctions without preserving sheaf structure may appear intelligent while undermining worldhood.

Taken together, these results establish a unified ontology in which physics, cognition, and technology are governed by the same principle: existence as the irreversible expenditure of possibility. Contemporary artificial intelligence systems derive their power from avoiding this expenditure, externalizing constraint in order to preserve reversibility. This choice explains both their remarkable performance and their characteristic shallowness.

Intelligence, on the account developed here, is not the capacity to compute indefinitely, but the capacity to live with the consequences of having acted. Worldhood is not an optional metaphysical embellishment, but an architectural necessity. Where irreversible commitment is systematically avoided, no world can arise. Where a system binds itself to its own history, meaning becomes possible.

\newpage
\appendix

\section{Appendix A: Formal Foundations of Irreversibility}

This appendix makes explicit several structural assumptions used throughout the paper concerning reversibility, history, and constraint.

\begin{definition}[Reversible Transition System]
Let $(X, T)$ be a transition system with state space $X$ and admissible transitions $T$. The system is reversible if for every $t \in T$ there exists an inverse $t^{-1} \in T$ such that
\[
t^{-1}(t(x)) = x \quad \forall x \in X.
\]
\end{definition}

\begin{definition}[Internal Irreversibility]
A system exhibits internal irreversibility if there exists a transition $t$ such that no admissible inverse exists within the system’s own transition set.
\end{definition}

\begin{lemma}[History Trivialization]
If a system is reversible, then any finite history $H$ is equivalent to the empty history with respect to future actional possibility.
\end{lemma}

\begin{proof}
By reversibility, every transition in $H$ admits an inverse. Composing inverses restores the initial state without loss. Hence $H$ induces no persistent constraint.
\end{proof}

This lemma formalizes the intuition underlying the No-World Theorem: reversibility trivializes history.

\section{Appendix B: Order-Theoretic Properties of Possibility Lattices}

We collect here several properties of the possibility lattices introduced in Sections 5 and 6.

\begin{definition}[Lower Set]
A subset $L \subseteq \mathcal{P}$ of a partially ordered set $(\mathcal{P}, \leq)$ is a lower set if
\[
x \in L \;\text{and}\; y \leq x \;\Rightarrow\; y \in L.
\]
\end{definition}

\begin{lemma}[Event Closure]
If an event maps $\mathcal{P}_t$ to a lower set $\mathcal{P}_{t+1}$, then no subsequent event can restore eliminated upper elements without introducing new structure.
\end{lemma}

\begin{proof}
Restoration would require reintroducing elements not in the lower set, violating monotonicity of the partial order. Hence restoration requires external augmentation.
\end{proof}

\begin{proposition}[Identity as Lattice Intersection]
The identity of a system over time is given by
\[
\bigcap_{t} \mathcal{P}_t.
\]
\end{proposition}

This proposition underwrites the process view of objecthood: identity is not a primitive, but the residue of constraint survival.

\section{Appendix C: RSVP Field Consistency Conditions}

We state here several consistency conditions required for the RSVP field equations to support world-bearing regimes.

\begin{definition}[Structured Entropy Growth]
Entropy growth is structured if
\[
\nabla S \cdot \mathbf{v} \neq 0.
\]
\end{definition}

\begin{lemma}[World Persistence Condition]
A necessary condition for persistent worldhood is
\[
\int_M \Phi |\mathbf{v}|^2 \, dx > 0
\]
over extended time intervals.
\end{lemma}

\begin{proof}
If the integral vanishes, either $\Phi = 0$ almost everywhere or $\mathbf{v} = 0$. In both cases, no directed expenditure of optionality occurs and world structure decays.
\end{proof}

\begin{remark}
This condition distinguishes living, historical regimes from equilibrium thermodynamic systems, which may possess entropy but lack directional commitment.
\end{remark}

\section{Appendix D: Categorical Structure of Care Operators}

We formalize here the categorical properties of care operators introduced in Section 8.

Let $\mathbf{Sh}(X)$ denote the category of sheaves over a base space $X$.

\begin{definition}[Care Functor]
A care operator is an endofunctor
\[
\mathcal{C} : \mathbf{Sh}(X) \rightarrow \mathbf{Sh}(X)
\]
that preserves finite limits.
\end{definition}

\begin{lemma}[Worldhood Preservation]
A care functor preserves worldhood if and only if it preserves pullbacks.
\end{lemma}

\begin{proof}
Pullbacks encode consistency across overlapping contexts. Failure to preserve them corresponds to collapse of global coherence. Preservation ensures that compressed representations remain historically meaningful.
\end{proof}

\begin{proposition}[Collapse as Limit Destruction]
A collapse operator corresponds to a functor that fails to preserve limits in $\mathbf{Sh}(X)$.
\end{proposition}

This proposition formalizes the distinction between attentional compression and semantic erasure.

\begin{thebibliography}{99}

\bibitem{heidegger1927}
Heidegger, Martin. 1927.
\textit{Being and Time}.
Translated by John Macquarrie and Edward Robinson.
New York: Harper \& Row, 1962.

\bibitem{heidegger1954}
Heidegger, Martin. 1954.
\textit{The Question Concerning Technology}.
In \textit{The Question Concerning Technology and Other Essays},
translated by William Lovitt, 3--35.
New York: Harper \& Row, 1977.

\bibitem{heidegger1938}
Heidegger, Martin. 1938.
\textit{The Age of the World Picture}.
In \textit{The Question Concerning Technology and Other Essays},
translated by William Lovitt, 115--154.
New York: Harper \& Row, 1977.

\bibitem{thomson2025}
Thomson, Iain D. 2025.
\textit{Heidegger on Technology’s Danger and Promise in the Age of AI}.
Cambridge: Cambridge University Press.
Published online 26 February 2025.

\bibitem{krouglov2025}
Krouglov, Alexander Yu. 2025.
Review of \textit{Heidegger on Technology’s Danger and Promise in the Age of AI},
by Iain D. Thomson.
\textit{Social Epistemology Review and Reply Collective} 14 (4): 13--14.

\bibitem{dreyfus1992}
Dreyfus, Hubert L. 1992.
\textit{What Computers Still Can’t Do: A Critique of Artificial Reason}.
Cambridge, MA: MIT Press.

\bibitem{dreyfus2007}
Dreyfus, Hubert L. 2007.
Why Heideggerian AI Failed and How Fixing It Would Require Making It More Heideggerian.
\textit{Philosophical Psychology} 20 (2): 247--268.

\bibitem{newell_simon1976}
Newell, Allen, and Herbert A. Simon. 1976.
Computer Science as Empirical Inquiry: Symbols and Search.
\textit{Communications of the ACM} 19 (3): 113--126.

\bibitem{dennett1987}
Dennett, Daniel C. 1987.
\textit{The Intentional Stance}.
Cambridge, MA: MIT Press.

\bibitem{searle1980}
Searle, John R. 1980.
Minds, Brains, and Programs.
\textit{Behavioral and Brain Sciences} 3 (3): 417--457.

\bibitem{ashby1956}
Ashby, W. Ross. 1956.
\textit{An Introduction to Cybernetics}.
London: Chapman \& Hall.

\bibitem{rosen1985}
Rosen, Robert. 1985.
\textit{Anticipatory Systems}.
Oxford: Pergamon Press.

\bibitem{friston2010}
Friston, Karl. 2010.
The Free-Energy Principle: A Unified Brain Theory?
\textit{Nature Reviews Neuroscience} 11: 127--138.

\bibitem{barandes2023}
Barandes, Jacob. 2023.
\textit{Unistochastic Quantum Theory}.
Oxford: Oxford University Press.

\bibitem{noble1984}
Noble, David F. 1984.
\textit{Forces of Production: A Social History of Industrial Automation}.
New York: Oxford University Press.

\end{thebibliography}

\end{document}
